\documentclass[11pt, a4paper]{article}
\usepackage[a4paper, margin=0.85in]{geometry}
\usepackage{mathptmx}
\usepackage{amsmath, amssymb}
\usepackage{graphicx}
\usepackage{booktabs}
\usepackage{hyperref}
\usepackage{caption}
\usepackage{subcaption}
\usepackage{enumitem}
\usepackage{float}
\usepackage{microtype}
\usepackage{xcolor}

\hypersetup{
    colorlinks=true,
    linkcolor=blue!70!black,
    citecolor=blue!70!black,
    urlcolor=blue!70!black
}

\title{\textbf{Experiment 5: Explainable AI (XAI) with SHAP \& LIME and Algorithmic Fairness Auditing with Fairlearn}}
\author{\textbf{Course:} Applied Data Science (ADS) \quad | \quad \textbf{Domain:} Responsible Machine Learning}
\date{\textbf{Benchmark Dataset:} Adult Census Income ($N=32,561$) \quad | \quad \textbf{Date:} September 2026}

\begin{document}

\maketitle

\begin{abstract}
As machine learning models are increasingly deployed in high-stakes socioeconomic environments, model opacity and algorithmic bias present critical technical, ethical, and legal challenges. This report presents an end-to-end investigation of Explainable AI (XAI) and Algorithmic Fairness on the Adult Census Income benchmark dataset ($N=32,561$, 12 features). First, we train high-performance baseline models, selecting a LightGBM champion achieving \textbf{87.69\% accuracy} and \textbf{0.9310 ROC-AUC}. Second, we apply game-theoretic \textbf{SHAP (Shapley Additive exPlanations)} via \texttt{TreeExplainer} to derive globally consistent feature importance, directional beeswarm distributions, non-linear interaction dependencies, and local sample attributions. Third, we deploy \textbf{LIME (\texttt{lime.lime\_tabular})} to construct locally faithful linear surrogate explanations across distinct customer profiles, assessing cross-methodological concordance. Fourth, we perform an algorithmic bias audit using \textbf{Fairlearn}, uncovering substantial demographic disparities across \textit{Sex} (Demographic Parity Difference = 0.1714, Disparate Impact Ratio = 0.346) and \textit{Race}. Finally, we systematically implement and benchmark a tri-modal bias mitigation suite spanning pre-processing (sample reweighting), in-processing (\texttt{ExponentiatedGradient} with Equalized Odds), and post-processing (\texttt{ThresholdOptimizer}). Post-processing successfully reduced the Equalized Odds Difference by \textbf{35.6\%} to 0.0393, while pre-processing reduced Demographic Parity Difference by \textbf{49.4\%} to 0.0867, illustrating the fundamental Fairness-Performance Pareto frontier.
\end{abstract}

\vspace{0.5em}
\hrule
\vspace{1em}

\section{Aim \& Objectives}
\textbf{Aim:} To apply explainable AI (XAI) methods (SHAP \& LIME) for interpreting model predictions and evaluate algorithmic fairness using Fairlearn.

\noindent \textbf{Objectives:}
\begin{enumerate}[leftmargin=2em]
    \item Implement global (SHAP summary bar and beeswarm distributions) and local (SHAP waterfall and LIME tabular) explanation frameworks.
    \item Audit demographic bias in machine learning models across protected sensitive attributes (\textit{Sex} and \textit{Race}) using Fairlearn's \texttt{MetricFrame}.
    \item Quantify disparity metrics including Selection Rate, Demographic Parity Difference, Equalized Odds Difference, and Disparate Impact Ratio.
    \item Propose, implement, and benchmark tri-modal bias mitigation strategies: Pre-processing (reweighting), In-processing (\texttt{ExponentiatedGradient}), and Post-processing (\texttt{ThresholdOptimizer}).
    \item Characterize the empirical Pareto trade-off between predictive accuracy and algorithmic fairness.
\end{enumerate}

\section{Theoretical Foundations of XAI \& Algorithmic Fairness}

\subsection{SHAP (Shapley Additive exPlanations)}
Rooted in cooperative game theory, SHAP allocates the payoff of an ensemble prediction among individual input features according to their marginal contributions across all possible feature coalitions $S \subseteq F \setminus \{i\}$:
\begin{equation}
\phi_i(v) = \sum_{S \subseteq F \setminus \{i\}} \frac{|S|!(|F| - |S| - 1)!}{|F|!} \left[ v(S \cup \{i\}) - v(S) \right]
\end{equation}
SHAP is uniquely characterized by four fundamental mathematical properties:
\begin{itemize}[leftmargin=1.5em]
    \item \textbf{Efficiency:} The sum of Shapley values equals the difference between model output and base expected value: $\sum_{i} \phi_i = f(x) - \mathbb{E}[f(X)]$.
    \item \textbf{Symmetry:} If features $i$ and $j$ contribute equally to all coalitions, $\phi_i = \phi_j$.
    \item \textbf{Dummy / Null Player:} If feature $i$ contributes nothing to any coalition, $\phi_i = 0$.
    \item \textbf{Additivity:} For ensemble models $f = f_1 + f_2$, the attributions satisfy $\phi_i(f) = \phi_i(f_1) + \phi_i(f_2)$.
\end{itemize}

\subsection{LIME (Local Interpretable Model-agnostic Explanations)}
LIME explains individual predictions by fitting an interpretable surrogate model $g \in G$ (e.g., sparse ridge regression) over localized perturbations $z'$ weighted by an exponential distance kernel $\pi_x(z)$:
\begin{equation}
\xi(x) = \arg\min_{g \in G} \mathcal{L}(f, g, \pi_x) + \Omega(g), \quad \text{where } \pi_x(z) = \exp\left(-\frac{D(x,z)^2}{\sigma^2}\right)
\end{equation}
While SHAP provides axiomatic guarantees, LIME offers model-agnostic operational flexibility.

\subsection{Algorithmic Fairness Formulations}
Let $Y \in \{0, 1\}$ denote the true ground truth label, $\hat{Y} \in \{0, 1\}$ denote the model prediction, and $A \in \{a, b\}$ denote the sensitive protected attribute:
\begin{itemize}[leftmargin=1.5em]
    \item \textbf{Demographic Parity (Statistical Parity):} Requires equal positive prediction rates across groups:
    \begin{equation}
    P(\hat{Y}=1 \mid A=a) = P(\hat{Y}=1 \mid A=b)
    \end{equation}
    \item \textbf{Equalized Odds:} Requires equal True Positive Rates (TPR) and False Positive Rates (FPR) across groups:
    \begin{equation}
    P(\hat{Y}=1 \mid Y=y, A=a) = P(\hat{Y}=1 \mid Y=y, A=b) \quad \forall y \in \{0, 1\}
    \end{equation}
    \item \textbf{Disparate Impact Ratio (80\% Four-Fifths Rule):} The selection rate ratio must satisfy:
    \begin{equation}
    \text{DIR} = \frac{\min(P(\hat{Y}=1 \mid A=a), P(\hat{Y}=1 \mid A=b))}{\max(P(\hat{Y}=1 \mid A=a), P(\hat{Y}=1 \mid A=b))} \ge 0.80
    \end{equation}
\end{itemize}

\section{Experimental Methodology \& Benchmark Results}

\begin{table}[H]
\centering
\caption{Comprehensive Benchmark: Performance vs. Fairness Disparity across Mitigation Strategies}
\label{tab:mitigation_benchmark}
\small
\begin{tabular}{lcccccc}
\toprule
\textbf{Model Pipeline} & \textbf{Acc} & \textbf{Bal Acc} & \textbf{F1} & \textbf{DP Diff (Sex)} & \textbf{EO Diff (Sex)} & \textbf{DIR (Sex)} \\
\midrule
1. Baseline Unmitigated (LightGBM) & 87.69\% & 80.65\% & 0.7240 & 0.1714 & 0.0611 & 0.3464 \\
2. Pre-Processing (Reweighted) & 86.92\% & 78.71\% & 0.6983 & 0.0867 & 0.2029 & 0.6086 \\
3. In-Processing (ExponentiatedGradient) & 85.69\% & 75.72\% & 0.6553 & 0.0973 & 0.0656 & 0.5292 \\
4. Post-Processing (ThresholdOptimizer) & 85.78\% & 77.14\% & 0.6719 & 0.1085 & 0.0393 & 0.5253 \\
\bottomrule
\end{tabular}
\end{table}

\begin{figure}[H]
    \centering
    \begin{subfigure}[b]{0.48\textwidth}
        \includegraphics[width=\textwidth]{plots/exp5_shap_summary_bar.png}
        \caption{Global Mean Absolute SHAP Ranking}
    \end{subfigure}
    \hfill
    \begin{subfigure}[b]{0.48\textwidth}
        \includegraphics[width=\textwidth]{plots/exp5_shap_beeswarm.png}
        \caption{SHAP Beeswarm Value Dispersion}
    \end{subfigure}
    \caption{SHAP Global Interpretability Suite on Validation Cohort ($N=800$).}
    \label{fig:shap_global}
\end{figure}

\begin{figure}[H]
    \centering
    \begin{subfigure}[b]{0.48\textwidth}
        \includegraphics[width=\textwidth]{plots/exp5_shap_dependence.png}
        \caption{SHAP Dependence: Age vs. Hours per Week}
    \end{subfigure}
    \hfill
    \begin{subfigure}[b]{0.48\textwidth}
        \includegraphics[width=\textwidth]{plots/exp5_shap_waterfall.png}
        \caption{Local SHAP Waterfall Attribution (Sample \#0)}
    \end{subfigure}
    \caption{SHAP Non-Linear Feature Interactions and Local Sample Attributions.}
    \label{fig:shap_local}
\end{figure}

\begin{figure}[H]
    \centering
    \begin{subfigure}[b]{0.48\textwidth}
        \includegraphics[width=\textwidth]{plots/exp5_lime_local_explanations.png}
        \caption{LIME Local Surrogate Decision Rules}
    \end{subfigure}
    \hfill
    \begin{subfigure}[b]{0.48\textwidth}
        \includegraphics[width=\textwidth]{plots/exp5_xai_shap_vs_lime_comparison.png}
        \caption{Cross-XAI Concordance (SHAP vs. LIME)}
    \end{subfigure}
    \caption{LIME Tabular Explanations and Cross-Method Concordance Audit.}
    \label{fig:lime_explanations}
\end{figure}

\begin{figure}[H]
    \centering
    \begin{subfigure}[b]{0.48\textwidth}
        \includegraphics[width=\textwidth]{plots/exp5_fairness_audit_disparity.png}
        \caption{Fairlearn Baseline Disparities (Sex \& Race)}
    \end{subfigure}
    \hfill
    \begin{subfigure}[b]{0.48\textwidth}
        \includegraphics[width=\textwidth]{plots/exp5_fairness_mitigation_tradeoff.png}
        \caption{Fairness-Accuracy Pareto Frontier}
    \end{subfigure}
    \caption{Algorithmic Fairness Audit and Empirical Pareto Trade-off Analysis.}
    \label{fig:fairness_analysis}
\end{figure}

\section{Discussion \& Key Findings}
\begin{enumerate}[leftmargin=1.5em]
    \item \textbf{Global Feature Dominance:} SHAP analysis established \texttt{Capital\_Gain}, \texttt{Age}, \texttt{Education\_Num}, and \texttt{Relationship} as the top drivers of high-income predictions. Marital status and relationship status act as strong socio-demographic proxies.
    \item \textbf{Cross-XAI Concordance:} For Sample \#0, both SHAP and LIME identified identical top positive drivers (\texttt{Education\_Num}, \texttt{Hours\_per\_week}) with consistent directional signs, validating local surrogate fidelity.
    \item \textbf{Baseline Fairness Deficiency:} The unmitigated champion model exhibited severe demographic disparity: male selection rate was 26.2\% compared to only 9.1\% for females ($\text{DIR} = 0.346$, violating the EEOC 80\% rule).
    \item \textbf{Mitigation Effectiveness:} Pre-processing reweighting achieved the greatest reduction in Demographic Parity Difference (0.0867 vs 0.1714, a 49.4\% improvement). Post-processing threshold optimization achieved the lowest Equalized Odds Difference (0.0393), proving that group-specific threshold calibration is superior for error-rate parity.
\end{enumerate}

\section{Conclusion}
Experiment 5 demonstrated the complete cycle of modern Responsible AI: extracting rigorous explanations through cooperative game theory and surrogate modeling, diagnosing systemic algorithmic bias, and strategically applying mitigation algorithms to achieve equitable outcomes on the Pareto efficiency frontier.

\end{document}
